% =====================================================================
% Section 4: Method — EvidenceFlow
% =====================================================================
%
% Status: Draft §4 Method v8 (2026-05-17)
%   v8: §4 opening rewritten as a Figure 1 walkthrough per advisor
%       feedback ("opening should be a brief description of the
%       method figure"). Body §4.1/§4.2/§4.3 untouched in this revision.
%   v7.5 (pre-opening): §4.3 reframed as Evidence-Aware Retrieval
%       Optimization, 4 equations, problem-then-solution framing.
%   v7.4: bodies stable, pre-opening checkpoint archived.
%
% Method name: \methodname (defined via \newcommand in main file)
% Forward reference: fig:overview (figure to be added in figures/)
%
% =====================================================================

\section{Method: \methodname{}}
\label{sec:method}

\begin{figure}[t]
\centering
\fbox{\parbox{0.92\linewidth}{\centering\vspace{1.5em}
\textit{Overview of \methodname{}.}\\[0.3em]
Indexing: corpus $\to$ fact-source substrate.\\
Retrieval: $G_q \to H_q \to \pi_q \to s_{\mathrm{ppr}}$.\\
Optimization: evidence-need mining + coverage-aware reranking $\to R_q$.
\vspace{1.5em}}}
\caption{Overview of \methodname{}. Indexing builds a fact-source
substrate that links each OpenIE fact to its source passage through
endpoint and source edges. Given a query, a query-local subgraph
$G_q$ is induced over this substrate, and a personalized PageRank on
the derived graph $H_q$ produces an evidence flow $\pi_q$ and a base
passage score $s_{\mathrm{ppr}}$. Two evidence-aware steps then
extend the candidate pool through evidence-need mining and
form the final ranked evidence list $R_q$ via coverage-aware
reranking.}
\label{fig:overview}
\end{figure}

Figure~\ref{fig:overview} gives an overview of \methodname{}. At
indexing time, the corpus is parsed into OpenIE facts whose head and
tail phrases are linked through two complementary edge types,
endpoint edges across passages and source edges within a passage,
yielding a fact-source substrate that is built once for $\mathcal{D}$
(\S\ref{sec:method:substrate}). At query time, the question first
selects a query-local subgraph $G_q$ over this substrate by combining
dense entry passages with the provenance of facts whose endpoints
match question-side anchors, after which a personalized PageRank on a
derived graph $H_q$ propagates mass through fact nodes and reads out
a passage score $s_{\mathrm{ppr}}$
(\S\ref{sec:method:retrieval}). The local ranking is then refined by
two evidence-aware steps over the same substrate: relation-grounded
expansion admits passages whose relation phrases match $q$ even when
they fall outside $G_q$, and coverage-aware reranking selects the
final list $R_q$ greedily under a noisy-OR coverage model over
question-side requirements (\S\ref{sec:method:enhancement}).

\subsection{Offline Indexing: A Fact-Source Substrate}
\label{sec:method:substrate}

Our indexing differs from the standard passage--phrase formulation
of HippoRAG-style systems in a single move: we promote each OpenIE
fact to a graph node and tie it to its source passage. The
resulting \emph{fact-source} substrate is indexed once for the
corpus, on which the propagation of \S\ref{sec:method:retrieval}
runs.

\paragraph{Knowledge Extraction.}
For each document $d \in \mathcal{D}$, a language model extracts
OpenIE facts $f=(h,r,t)$ with argument phrases $h, t$ and relation
phrase $r$
\citep{angeli2015leveraging,gutierrez2025hipporag}. We retain every
fact as a first-class indexing unit and bind it to its source via
$\mathrm{src}(f)=d$, yielding a global fact set $\mathcal{F}$ with
$\mathcal{F}_d = \{f\in\mathcal{F} : \mathrm{src}(f)=d\}$. Retaining
$f$ together with $\mathrm{src}(f)$ keeps the relational triple
intact, so a fact's role in retrieval is determined by both its
argument structure and its source passage rather than by its
endpoints alone.

\paragraph{Graph Construction.}
Beyond the standard passage--phrase edges of earlier OpenIE
graphs, the fact-source substrate adds a \emph{source edge} that
ties each fact to its provenance. \emph{Endpoint edges} connect
two facts that share at least one informative argument phrase
under $\{h,t\}\cap\{h',t'\}\neq\emptyset$ when the phrases come
from different source passages, encoding cross-document mentions
of the same entity. \emph{Source edges} connect facts that share
$\mathrm{src}(f)=\mathrm{src}(f')$, encoding co-assertion within a
passage. Each fact also carries a passage link to $\mathrm{src}(f)$,
so passage scores can later be read out by aggregating mass on
local facts. Dense embeddings of fact units and relation phrases
are pre-computed; the substrate is materialized as an adjacency
index with a hub-degree cap of $30$ and a synonymy threshold of
$0.85$ on dense endpoint similarity
(Appendix~\ref{app:hyperparams}). Endpoint edges thus carry mass
across passages while source edges carry mass within a passage.

\subsection{Evidence-Driven Local Graph Retrieval}
\label{sec:method:retrieval}

\paragraph{Query-local subgraph.}
Given $q$, \methodname{} first induces a query-local subgraph $G_q$
that selects the part of the offline fact-source substrate
considered for the query. We extract entity anchors
$A(q) \subseteq \mathcal{V}_{\mathrm{phr}}$ from $q$ using a
controlled language model that operates on the question only,
never on retrieved documents, and a standard dense retriever
provides an entry signal $\mathrm{Entry}(q) \subseteq \mathcal{D}$.
The seed passage set combines $\mathrm{Entry}(q)$ with the source
passages of facts whose endpoints overlap $A(q)$. From this seed
set, $G_q$ is the bounded-hop expansion that follows endpoint and
source edges over the substrate of \S\ref{sec:method:substrate}.
$G_q$ keeps the activated region compact while still admitting
passages reachable through multi-hop entity chains even when their
surface similarity to $q$ is low.

\paragraph{Local Graph Diffusion.}
The subgraph $G_q$ selects the activated region but does not itself
produce a passage ranking. We therefore build a derived graph
$H_q$ over the local fact nodes and run a personalized PageRank on
it. Propagation is materialized over OpenIE \emph{facts}, so it
traverses multi-hop chains at the relational unit level rather
than over phrase nodes; endpoint edges carry mass across passages
while source edges carry mass within a passage. The resulting
query-conditioned mass distribution $\pi_q$, which we call the
\emph{evidence flow}, is therefore both local to $q$ and grounded
in the offline graph. A fact $f$ enters $\mathcal{F}_q$ when
$\mathrm{src}(f)$ is a passage in $G_q$, so the phrase side enters
$H_q$ only through endpoint edges. We solve the local PPR equation
\begin{equation}
\label{eq:local-ppr}
\pi_q = \alpha\, s_q + (1-\alpha)\, P_q^{\!\top}\, \pi_q,
\end{equation}
where $s_q$ places mass on facts whose arguments include an anchor
in $A(q)$, weighted by entry scores of the corresponding passages,
and $P_q$ is the row-normalized transition matrix over the local
fact edges, with no edge-type balancing beyond the hub cap of
\S\ref{sec:method:substrate}. We use local push with reset weight
$\alpha=0.2$ and residual threshold $10^{-6}$. The passage score
for $d$ aggregates stationary mass on local facts grounded in $d$,
$s_{\mathrm{ppr}}(d \mid q) = |\mathcal{F}_{q,d}|^{-1/2}
\sum_{f\in\mathcal{F}_{q,d}}\pi_q(f)$, with
$\mathcal{F}_{q,d}=\{f\in\mathcal{F}_q:\mathrm{src}(f)=d\}$. The
square-root normalization is a middle ground between summing fact
mass and averaging it: it prevents passages with many extracted
facts from dominating only by length, while preserving evidence
from multiple activated facts. Because propagation is restricted
to $H_q$, $s_{\mathrm{ppr}}$ rewards passages connected to $q$
through entity-relation evidence rather than passages independently
similar to $q$. We refer to this stage as \emph{local graph
retrieval} when discussing ablations
(\S\ref{sec:experiments:ablation}).

\subsection{Evidence-Coverage-Aware Retrieval Optimization}
\label{sec:method:enhancement}

The local-graph-retrieval ranking $s_{\mathrm{ppr}}$ is graph-aware
but treats each passage in isolation, so two reader-facing problems
remain. A passage on a multi-hop chain can fall outside $G_q$ and
receive little $\pi_q$-mass even when its relation phrase exactly
matches what $q$ asks for. The head of the ranking can also crowd
around a single entity neighborhood, leaving other entities or
relations demanded by $q$ uncovered in any short prefix. We
extend $s_{\mathrm{ppr}}$ with two evidence-aware steps. We first
\emph{mine the evidence the query still needs}, enlarging the
candidate pool through query-relevant relations. We then apply a
set-based reranking step to the resulting candidate set, so that
the final reader prefix addresses distinct question-side
requirements rather than repeating the strongest neighborhood.

\paragraph{Evidence-need mining.}
This step targets the bridge-underweighting problem by promoting
passages whose relation phrase carries query-relevant evidence,
including passages outside $G_q$. For a candidate fact
$e=(h_e, r_e, t_e)$ in $\mathcal{F}$ with $\mathrm{src}(e)=d$, we
score $d$ by combining the largest $\pi_q$-mass on any local fact
$f=(h_f, r_f, t_f) \in \mathcal{F}_q$ sharing an argument phrase
with $e$ and the dense compatibility between $r_e$ and $q$,
\begin{equation}
\label{eq:expansion-score}
\begin{aligned}
s_{\mathrm{exp}}(d \mid q)
= {}& \max_{\substack{e:\,\mathrm{src}(e)=d \\ f \in \mathcal{F}_q}}
\Bigl\{\\[-0.3em]
&\pi_q(f)\,\bigl[\sigma(\mathbf{r}_e,\mathbf{q})\bigr]_{+}\\
&{}\cdot\mathbf{1}\!\left[
\{h_f,t_f\}\cap\{h_e,t_e\}\neq\emptyset
\right]\Bigr\},
\end{aligned}
\end{equation}
where $\sigma$ is cosine similarity, $\mathbf{r}_e, \mathbf{q}$ are
dense embeddings of the relation phrase and the question, and
$[x]_{+}=\max(x,0)$ clips negative similarities. The max operator
lets a single high-confidence relation bridge rescue a passage and
avoids score inflation from many low-confidence matches in the
same neighborhood. The base score linearly combines the two
signals,
\begin{equation}
\label{eq:base-score}
s_{\mathrm{base}}(d \mid q)
=
s_{\mathrm{ppr}}(d \mid q)
+
\eta\,s_{\mathrm{exp}}(d \mid q),
\end{equation}
with $\eta$ tuned once on a development split. Passages without an
evidence-need path have $s_{\mathrm{exp}}(d \mid q)=0$, so the base
score reduces to local-graph retrieval in the absence of
relational evidence. Let $C_q$ denote the union of passages scored
by local-graph retrieval and passages admitted by evidence-need
mining.

\paragraph{Coverage-aware reranking.}
This step targets the entity-neighborhood crowding problem by
selecting passages that collectively address the distinct
question-side requirements $B(q)=\{b_1,b_2,\ldots\}$ extracted from
$q$ by the same controlled language model used for $A(q)$. For a
passage $d$ and requirement $b$, $\phi_{d,b}\in[0,1]$ is computed
from NV-Embed-v2 cosine similarity between the passage text and
the requirement subquery, min--max normalized over $C_q$. When $b$
depends on a resolved bridge entity, $\phi_{d,b}$ is multiplied by
the normalized compatibility between $d$ and that entity;
otherwise the multiplier is $1$. Defining the noisy-OR coverage of
$R$ as
$\mathrm{Cov}_q(R)=\sum_{b\in B(q)}\bigl[1-\prod_{d'\in R}(1-\phi_{d',b})\bigr]$,
the marginal coverage gain of adding $d$ is
\begin{equation}
\label{eq:marginal-gain}
\Delta_q(d \mid R) = \mathrm{Cov}_q(R\cup\{d\}) - \mathrm{Cov}_q(R).
\end{equation}
We then form $R_q$ by greedy MMR-style set selection
\citep{carbonell1998mmr,lin2011submodular}, at each step adding
\begin{equation}
\label{eq:greedy-rerank}
d^{*}
=
\operatorname*{arg\,max}_{d \in C_q \setminus R}
\lambda\,s_{\mathrm{base}}(d \mid q)
+
(1-\lambda)\,\Delta_q(d \mid R),
\end{equation}
with $\lambda \in [0,1]$ tuned once on a development split. Because
per-requirement coverage saturates, passages duplicating
already-covered neighborhoods receive diminishing marginal gain,
while passages addressing previously uncovered requirements are
promoted. The reader consumes any prefix of $R_q$.
